\documentclass[11pt]{article}
\usepackage[utf8]{inputenc}
\usepackage{amsmath,amssymb}
\usepackage{graphicx}
\usepackage{booktabs}
\usepackage{authblk}
\usepackage[margin=1in]{geometry}
\usepackage[hidelinks]{hyperref}
\usepackage{xurl}
\usepackage{setspace}

\title{The Spinor Residue:\\
$g=2$ and Stern--Gerlach without Relativity}
\author[1]{Claes Johnson}
\affil[1]{KTH Royal Institute of Technology, SE-100 44 Stockholm, Sweden \\
\texttt{claesjohnson@gmail.com}}
\date{\today \\[0.4em]
\small\textit{Argument by the author; drafting prepared with AI assistance (Claude, Anthropic), reviewed and verified by the author.}}

\begin{document}
\maketitle

\begin{abstract}
The charges-as-configurations reading of RealQM removes relativity and gives the electron's magnetic
moment as a quantised current-winding: an orbital winding $e^{im\phi}$ carries $\mu_z = -m\,\mu_B$ with
gyromagnetic factor $g=1$. This leaves exactly one residue --- the electron's \emph{spin} moment has
$g=2$, twice the orbital value per unit angular momentum --- and one experiment that RealQM in its present
form cannot reproduce: \emph{Stern--Gerlach}, whose two-spot splitting is the signature of two-valued
spin $\tfrac12$. We make the residue precise, and argue that it is \emph{not} a demand for relativity. By
Lévy-Leblond's theorem $g=2$ follows from a first-order (linearised), \emph{non-relativistic} spinor wave
equation; the ``2'' is the SU(2) double cover of the rotation group, a geometric fact, not a consequence
of Lorentz invariance. Concretely, writing the kinetic energy with $\boldsymbol\sigma\cdot\mathbf p$
rather than $p^2$ produces, in a field, the term $-q\,\boldsymbol\sigma\cdot\mathbf B$ --- the $g=2$
Zeeman coupling --- automatically. The target for RealQM is therefore identified without importing
relativity: equip the non-overlapping domains with the $\pm\tfrac12$ spinor label the base theory already
suggests, implemented as a $\boldsymbol\sigma\cdot\mathbf p$ structure, so that the moment becomes
two-valued ($\mu_z=\pm\mu_B$) and Stern--Gerlach's two spots and $g=2$ follow together. We are explicit
that this is a \emph{target}, not a completed derivation: what is established is that the residue is
geometric and reachable, not relativistic.
\end{abstract}

\setstretch{1.12}

\section{The one residue}
\label{sec:residue}

In the current-bearing form of RealQM the magnetic moment is that of the charge current,
$\boldsymbol\mu=\tfrac12\int\mathbf r\times\mathbf J$, and an orbital phase winding $e^{im\phi}$ gives the
quantised value $\mu_z=-m\,\mu_B$~\cite{ComplexRealQM,NuclearMoment}. This is gyromagnetic factor
$g=1$: one unit of orbital angular momentum, $L_z=m\hbar$, yields $m$ Bohr magnetons. The
charges-as-configurations picture~\cite{Charges} builds on this to remove any need for relativity: sizes
and moments are properties of Coulomb configurations, a ``free'' electron is only a weakly-confined one,
and the winding it can afford gives $\mu_B$ with no rest-energy scale invoked.

Exactly one thing is left unaccounted. The electron's \emph{spin} moment is not $g=1$ but
\textbf{$g=2$}: a spin of only $\hbar/2$ produces a \emph{full} Bohr magneton, twice what an orbital
current of the same angular momentum would give. An orbital winding cannot supply this factor of two; it
is the signature of spin, and RealQM has, so far, no spin-$\tfrac12$ moment. This is the residue.

\section{Its experimental face: Stern--Gerlach}
\label{sec:sg}

The residue is not academic; it is the whole content of the Stern--Gerlach experiment. A beam of silver
atoms --- each with a single valence electron in an $L=0$ (orbitally spherical) ground state --- passes
through an inhomogeneous magnetic field and splits into \emph{two} discrete spots. Two spots is the
signature of a two-valued moment, $\mu_z=\pm\mu_B$, i.e.\ spin $\tfrac12$ with $g=2$; an orbital moment
would split into an \emph{odd} number $2l+1$, and for the $L=0$ silver electron would not split at all
($m=0$, $\mu_z=0$).

RealQM in its present form therefore \emph{fails} Stern--Gerlach: with only orbital windings the $L=0$
valence electron carries no moment and the beam would pass undeflected. Reproducing the two-spot pattern
requires precisely the two-valued spin-$\tfrac12$ structure that is missing. Stern--Gerlach and $g=2$ are
one problem: the two spots \emph{are} $g=2$ made visible.

\section{The residue is geometric, not relativistic}
\label{sec:geom}

It is tempting to conclude that spin, $g=2$, and Stern--Gerlach force relativity into the theory, since in
the textbook account $g=2$ falls out of the Dirac equation. This conclusion is false, and the point is
decisive for the present programme.

\paragraph{Lévy-Leblond.} In 1967 Lévy-Leblond showed that $g=2$ follows from a first-order (linearised)
\emph{non-relativistic}, Galilean spinor wave equation~\cite{LevyLeblond}. Relativity is sufficient to
produce $g=2$ but not necessary; the factor of two survives the non-relativistic limit because it does not
come from Lorentz invariance in the first place.

\paragraph{The mechanism.} Write the kinetic energy not as $p^2/2m$ but as
$(\boldsymbol\sigma\cdot\mathbf p)^2/2m$ --- legal for free motion since
$(\boldsymbol\sigma\cdot\mathbf p)^2 = p^2$. In a magnetic field, minimal coupling
$\mathbf p\to\mathbf p-q\mathbf A$ gives
\begin{equation}
\frac{(\boldsymbol\sigma\cdot(\mathbf p-q\mathbf A))^2}{2m}
= \frac{(\mathbf p-q\mathbf A)^2}{2m} \;-\; \frac{q}{2m}\,\boldsymbol\sigma\cdot\mathbf B ,
\label{eq:pauli}
\end{equation}
and the last term is exactly the $g=2$ spin--Zeeman coupling, appearing \emph{automatically} the moment
the dynamics is linearised in $\boldsymbol\sigma\cdot\mathbf p$. No relativity is used in
\eqref{eq:pauli}; it is the ordinary Pauli Hamiltonian.

\paragraph{The geometry.} The ``2'' is the \textbf{SU(2) double cover of the rotation group}: a spinor
returns to itself only after $720^\circ$, rotating at half angle. That half is reciprocally the double in
the moment. It is the same species of fact as the winding-number topology that fixes $\mu_z=-m\mu_B$ in
the orbital case --- a statement about how spin-$\tfrac12$ objects transform, geometric and
scale-free, owing nothing to rest energy.

\section{The target for RealQM}
\label{sec:target}

The residue thus lands on the reachable side of the ledger. What RealQM needs is not relativity but the
right \emph{spinor structure} on its domains, and the base theory already points at it: the $\psi_i$ on
$\Omega_i$ acquiring an internal $\pm\tfrac12$ label. The programme is to implement that label not as a
bolt-on but as a $\boldsymbol\sigma\cdot\mathbf p$ linearisation of the one-electron dynamics on each
domain, so that:
\begin{itemize}
\item the spin--Zeeman term $-\tfrac{q}{2m}\boldsymbol\sigma\cdot\mathbf B$ of \eqref{eq:pauli} appears,
giving $g=2$;
\item the moment becomes \emph{two-valued}, $\mu_z=\pm\mu_B$, so that an $L=0$ electron in an
inhomogeneous field splits into \emph{two} beams --- Stern--Gerlach --- rather than none;
\item the non-overlap/free-boundary structure, and the current $\mathbf J$ of the complex form, are
retained unchanged, since \eqref{eq:pauli} is a rewriting of the kinetic term, not a new interaction.
\end{itemize}

We state plainly that this is a \emph{target}, not a finished calculation: we have not carried the
$\boldsymbol\sigma\cdot\mathbf p$ structure through the multiphase free-boundary problem, and doing so ---
in particular checking that the two-valued moment coexists with the geometric realisation of Pauli
exclusion by non-overlap without double-counting the spin --- is the open foundational step. What is
established is the map: $g=2$ and Stern--Gerlach are one residue, that residue is the SU(2) spinor factor
of two, and it is obtainable non-relativistically, so it belongs inside RealQM rather than outside it.

\section{Conclusion}
\label{sec:concl}

The charges-as-configurations picture closes the free-charge and relativity problems and leaves a single,
sharp residue: the electron's spin factor $g=2$, whose experimental embodiment is the two-spot
Stern--Gerlach splitting. That residue is not a relativistic quantity smuggled in through the Dirac
equation; by Lévy-Leblond it is a non-relativistic, geometric consequence of the SU(2) spinor structure,
delivered by a $\boldsymbol\sigma\cdot\mathbf p$ linearisation. RealQM therefore has a concrete,
relativity-free target for spin: give its domains the $\pm\tfrac12$ label as a $\boldsymbol\sigma\cdot
\mathbf p$ structure, and $g=2$ and Stern--Gerlach's two beams should follow together. Until that is
carried through the multiphase problem the point is a target, not a result --- but it is the right size of
target, and it is inside the theory, not beyond it.

\begin{thebibliography}{9}
\bibitem{ComplexRealQM} C.~Johnson, \emph{Complex Time-Dependent RealQM: Currents, Radiation, and
Magnetism from the Real-Space Continuum}, 2026.
\bibitem{NuclearMoment} C.~Johnson, \emph{Magnetism as Circulating Charge: Why the Nuclear Electron
Carries No Magnetic Moment}, 2026.
\bibitem{Charges} C.~Johnson, \emph{Elementary Charges as Configurations, not Objects: Why RealQM Needs No
Relativity}, 2026.
\bibitem{LevyLeblond} J.-M.~Lévy-Leblond, \emph{Nonrelativistic particles and wave equations},
Commun. Math. Phys. \textbf{6}, 286 (1967).
\end{thebibliography}

\end{document}
